\begin{abstract}
Increasing workload can slow individual operations and increase repeated processing, allowing congestion to persist after a demand surge. We study which dependencies an executable process model must preserve to prepare an effective response from logs recorded before sustained severe overload. The method learns service and routing behavior, compares feasible interventions under specified stresses, and fixes an action and observable trigger before independent evaluation. In the serial reference, retaining load dependencies reduces backlog and intervention cost by 38.2\% relative to static planning and 35.4\% relative to predictive component selection. A new evaluation covers 288 training datasets across six structures with parallel operations, alternative routes, multiple returns and competing resources. Under routing feedback, the complete pipeline with automatically discovered structure reduces mean normalized cost by 13.5\% relative to predictive selection; the adjusted 97.5\% difference interval is $[-0.1353,-0.0343]$. Selecting a structure that preserves return behavior and observed route constraints also improves decisions over standard Inductive Miner. Benefits depend on the feedback mechanism and process. The results establish why response preparation requires both decision-relevant load dependencies and an executable structure that preserves the meaning of interventions.
\keywords{Process discovery \and Business process simulation \and Congestion \and Response preparation \and Load-dependent behavior}
\end{abstract}

\section{Introduction}
In a complex information system, load can change the work required to complete each case. Individual operations slow down, and repeated processing becomes more frequent. Both effects occupy resources for longer, causing further work to accumulate. A temporary surge can therefore leave the system congested after arrivals return to their usual rate. Preparing a response requires understanding these feedback mechanisms as well as capacity constraints.

Different actions change different parts of the mechanism. Adding a worker increases processing capacity; reducing rework decreases the work generated by each case. In a process-aware information system, engineers must choose a feasible action, specify an observable trigger and assess system-wide consequences before deployment. Parallel branches, repeated subprocesses and shared resources make those consequences depend on more than the local queue where a problem first appears.

Predictive accuracy under ordinary operation does not directly answer this decision question. A load dependency can contribute little to prediction within the observed range yet influence the ranking of actions during overload. Removing it can leave a model that describes familiar conditions adequately but prepares the wrong response. Conversely, retaining a weakly estimated dependency can introduce an incorrect extrapolation. The modeling decision must therefore be assessed through its effect on the prepared action.

Business process simulation and prescriptive monitoring establish the foundations for learning models from logs and evaluating interventions~\cite{simod,alarms,hybridpolicy,rlpolicies}; workload-dependent durations are already used in hybrid simulation~\cite{deepsimulator,hybridpolicy}. We examine three connected issues in this setting: retaining dependencies relevant to response choice, separating predictive quality from decision quality, and testing whether choices remain effective under structural complexity and discovery errors.

We compare behavioral component selection and control-flow structure while holding the action catalog and simulation search fixed. The serial reference isolates component selection; six designed process classes test preparation with supplied structure and train-only discovery. A feedback rule and a known-generator planner separate the value of the learned representation from that of searching over actions. The contribution connects these modeling choices to independent response outcomes.

\section{Related Work}
Queue mining uses event logs to estimate waiting times~\cite{queue}. Research on load-dependent service shows that workload can also change processing time itself~\cite{load}. Studies of metastable failures describe how temporary disturbances initiate self-sustaining work amplification~\cite{metastable}. We study these mechanisms through finite-horizon process congestion and the actions available to limit it.

Process discovery and data-driven simulation supply executable models of system behavior. Inductive Miner and PM4Py recover process structure~\cite{inductive,pm4py}; SIMOD discovers and calibrates simulation models~\cite{simod}; Prosimos models differentiated resource behavior~\cite{prosimos}; and AgentSimulator uses a resource-centric decomposition~\cite{agents}. DeepSimulator combines discovered structure with learned timing and intercase features~\cite{deepsimulator}. 

Decision mining learns routing rules from case attributes~\cite{decisionmining}. Lifecycle-aware discovery distinguishes concurrency from interleaving to improve performance measurement~\cite{lifecycle}. We use a restricted process-tree discovery procedure to isolate how control-flow and routing constraints affect the intervention selected under endogenous workload feedback.

The distinction between prediction error and decision quality is established in predict-then-optimize and decision-focused learning. Smart Predict-then-Optimize measures predictive errors through their downstream optimization loss~\cite{spo}; decision-focused learning trains predictive models against decision objectives~\cite{dflsurvey}. Our contribution is not a new decision-focused training loss. We investigate component retention in an executable process model fitted before a sustained crisis. Here workload is endogenous: an intervention changes queues, queues change service and repeated processing, and those changes alter future workload. Historical logs provide observations of ordinary execution, rather than labels for the costs of alternative crisis responses. The engineering problem is to preserve the behavioral dependencies needed to compare those responses while respecting control flow, synchronization and shared resources.

The closest policy-learning work includes workload-aware resource allocation. Meneghello et al.~\cite{hybridpolicy} combine PPO with RIMS hybrid simulation. Their duration model uses work in progress and resource occupation, routing uses gateway probabilities, and the evaluation includes congestion caused by resource restrictions. Middelhuis et al.~\cite{resourcepolicy} learn resource assignments through deep reinforcement learning and score-based value approximation, with queues and ongoing assignments as state. Branchi et al.~\cite{rlpolicies} learn activity policies from event logs using explicit Markov decision processes and offline reinforcement learning. Their comparison of fixed gateway probabilities with decisions based on case attributes and history demonstrates the importance of the evaluation environment's representation.

Prescriptive monitoring addresses complementary parts of response design. White-box policies provide configurable filtering and ranking, with parameter optimization under resource constraints~\cite{whitebox}; alarm-based systems optimize intervention timing~\cite{alarms}; and causal methods estimate intervention effects from historical treatment data~\cite{causal}. Multi-objective prescription balances organizational KPIs~\cite{multikpi}. We take action effects as engineering inputs and evaluate both backlog cost and sustained-overload incidence. Simulation-quality frameworks distinguish control-flow, temporal and congestion fidelity~\cite{validation}; our evaluation adds the realized consequence of the action selected by the model.
